\documentclass[../../main.tex]{subfiles}% 注意这里的文件路径不能用 ./main.tex ,否则用latexmk编译子文件会报错
\graphicspath{{\subfix{./image/}}{\subfix{../../image/}}} % 指定图片引用路径,后续可以直接使用图片文件名
% 如果图片存放在上级目录,则使用 ../../image/ 作为备用路径

% 例如:
% \begin{figure}[H]
% \centering
% \includegraphics[width=0.3\textwidth]{图.png}
% \caption{}
% \label{figure:图}
% \end{figure}
% 注意:上述\label{}一定要放在\caption{}之后,否则引用图片序号会只会显示??.

\begin{document}

\section{反常积分敛散性判别}

\begin{theorem}[Cauchy收敛准则]
广义积分 \(\int_{a}^{\infty} f(x) \mathrm{d}x\) 收敛等价于对任意 \(\varepsilon > 0\)，存在 \(A > a\) 使得任意 \(x_1,x_2 > A\) 都有 \(\left|\int_{x_1}^{x_2} f(t) \mathrm{d}t\right| < \varepsilon\)。 
\end{theorem}

\begin{theorem}
设在 $[a, +\infty)$ 上 $f \geqslant 0$, $f$ 在 $[a, +\infty)$ 的任何有界子区间上可积, 则无穷积分 $\int_a^{+\infty} f(x) \, \mathrm{d}x$ 收敛的充要条件是: 存在 $M > 0$ 使得对任何 $b > a$ 都有
\[
\int_a^b f(x) \, \mathrm{d}x < M,
\]
即 $F(b) = \int_a^b f(x) \, \mathrm{d}x$ 对于任何 $b$ 有界.
\end{theorem}

\begin{theorem}[比较判别法]
设 \( f \) 和 \( g \) 在 \([a, +\infty)\) 上有定义, 对任意 \( b > a \), \( f \) 和 \( g \) 在 \([a, b]\) 可积, 且对充分大的 \( x \), 成立不等式 \( 0 \leqslant f(x) \leqslant g(x) \). 若 \( \int_a^{+\infty} g(x) \, \mathrm{d}x \) 收敛, 则 \( \int_a^{+\infty} f(x) \, \mathrm{d}x \) 收敛.
\end{theorem}

\begin{theorem}[比较判别法极限形式]
如果 \( f \) 和 \( g \) 在 \([a, +\infty)\) 上有定义且非负, 并且对任意 \( b > a \), \( f \) 和 \( g \) 在 \([a, b]\) 上可积, \( \lim\limits_{x \to +\infty} \frac{f(x)}{g(x)} = k \), 那么有

(1) 若 \( 0 < k < +\infty \), 则 \( \int_a^{+\infty} f(x) \, \mathrm{d}x \) 与 \( \int_a^{+\infty} g(x) \, \mathrm{d}x \) 同敛散;

(2) 若 \( k = 0 \), 则当 \( \int_a^{+\infty} g(x) \, \mathrm{d}x \) 收敛时, \( \int_a^{+\infty} f(x) \, \mathrm{d}x \) 也收敛;

(3) 若 \( k = +\infty \), 则当 \( \int_a^{+\infty} g(x) \, \mathrm{d}x \) 发散时, \( \int_a^{+\infty} f(x) \, \mathrm{d}x \) 也发散.
\end{theorem}

\begin{theorem}[A-D判别法]\label{theorem:A-D判别法}
设 \(f(x),g(x)\) 在任何闭区间上黎曼可积，其$f,g$在$x=a$处都有界.
\begin{enumerate}
\item {\heiti Abel判别法}：若 \(\int_{a}^{+\infty} f(x) \mathrm{d}x\) 收敛，并且$g(x)$在$[a,+\infty)$上单调有界,则 \(\int_{a}^{\infty} f(x)g(x) \mathrm{d}x\) 收敛.

\item {\heiti Dirichlet判别法}：若 \(\int_{a}^{x} f(x) \mathrm{d}x\)在$[a,+\infty)$上有界，并且\(g(x)\)在\([a,+\infty)\)上单调,$\underset{x\rightarrow +\infty}{\lim}g\left( x \right) =0$,则 \(\int_{a}^{\infty} f(x)g(x) \mathrm{d}x\) 收敛.
\end{enumerate}
\end{theorem}
\begin{remark}
Dirichlet判别法要强于Abel判别法.因为可以由Dirichlet判别法直接推出Abel判别法.证明如下 :

设 \(f(x),g(x)\) 在任何闭区间上黎曼可积，其$f,g$在$x=a$处都有界.若 \(\int_{a}^{+\infty} f(x) \mathrm{d}x\) 收敛，并且$g(x)$在$[a,+\infty)$上单调有界,则$\lim\limits_{x\rightarrow +\infty}g\left( x \right) \triangleq A\in \mathbb{R} $, 令$h\left( x \right) =g\left( x \right) -A$, 则$\lim\limits_{x\rightarrow +\infty}h\left( x \right) =0$, 并且$h\left( x \right) $与$g\left( x \right) $有相同单调性. 由$\int_a^{\infty}{f\left( x \right)}\mathrm{d}x$收敛可知, $\int_a^x{f\left( t \right)}\mathrm{d}t$在$\left[ a,+\infty \right) $上必有界. 从而由Dirichlet判别法可知$\int_a^{\infty}{f\left( x \right) h\left( x \right)}\mathrm{d}x$收敛. 于是
\begin{align*}
\int_a^{\infty}{f\left( x \right) g\left( x \right)}\mathrm{d}x=\int_a^{\infty}{f\left( x \right) h\left( x \right)}\mathrm{d}x+A\int_a^{\infty}{f\left( x \right)}\mathrm{d}x<+\infty .
\end{align*}
\end{remark}

\begin{example}
设 \(f(x)\) 在 \([0,+\infty)\) 中非负且递减，证明：\(\int_{0}^{+\infty} f(x) \mathrm{d}x\)，\(\int_{0}^{+\infty} f(x)\sin^{2}x \mathrm{d}x\) 同敛散性. 
\end{example}
\begin{proof}
(i) 若 \(\int_0^\infty f(x) \, \mathrm{d}x < \infty\)，则由条件可知
\begin{align*}
f(x) \sin^2 x \leqslant f(x), \quad \forall x \in [0, +\infty).
\end{align*}
故由比较判别法可得 \(\int_0^\infty f(x) \sin^2 x \, \mathrm{d}x < \infty\)。

(ii) 若 \(\int_0^\infty f(x) \sin^2 x \, \mathrm{d}x < \infty\)，则由 \(f\) 非负递减，故 \(\lim_{x \to +\infty} f(x)\) 存在且 \(\lim_{x \to +\infty} f(x) \geqslant 0\)。若 \(\lim_{x \to +\infty} f(x) \triangleq a > 0\)，则存在 \(M > 0\)，使得
\begin{align}
f(x) \sin^2 x > \frac{a}{2} \sin^2 x, \quad \forall x \in [M, +\infty). \label{eq:100.31}
\end{align}
又因为
\begin{align*}
\int_0^\infty \sin^2 x \, \mathrm{d}x = \lim_{b \to +\infty} \int_0^b \frac{1 - \cos 2x}{2} \, \mathrm{d}x = \frac{1}{2} \lim_{b \to +\infty} \left( b - \frac{\sin 2b}{2} \right),
\end{align*}
而上式右边极限不存在，所以 \(\int_0^\infty \sin^2 x \, \mathrm{d}x\) 发散。从而结合 \eqref{eq:100.31} 式，由比较判别法可知 \(\int_0^\infty f(x) \sin^2 x \, \mathrm{d}x\) 发散，矛盾！故 \(\lim_{x \to +\infty} f(x) = 0\)。

注意到
\begin{align*}
\int_0^\infty f(x) \sin^2 x \, \mathrm{d}x = \frac{1}{2} \int_0^\infty f(x) (1 - \cos 2x) \, \mathrm{d}x < \infty.
\end{align*}
即 \(\int_0^\infty f(x) (1 - \cos 2x) \, \mathrm{d}x < \infty\)。考虑 \(\int_0^\infty f(x) \cos 2x \, \mathrm{d}x\)，注意到
\begin{align*}
\int_0^C \cos 2x \, \mathrm{d}x = \frac{\sin 2C}{2} < 1, \quad \forall C > 0.
\end{align*}
又由于 \(f(x)\) 在 \([0, +\infty)\) 上单调递减趋于 0，故由狄利克雷判别法可知 \(\int_0^\infty f(x) \cos 2x \, \mathrm{d}x < \infty\)。因此
\begin{align*}
\int_0^\infty f(x) \, \mathrm{d}x = \int_0^\infty f(x) (1 - \cos 2x) \, \mathrm{d}x + \int_0^\infty f(x) \cos 2x \, \mathrm{d}x < \infty.
\end{align*}

(iii)当\(\int_0^\infty f(x) \, \mathrm{d}x\) 或 \(\int_0^\infty f(x) \sin^2 x \, \mathrm{d}x\) 发散时,实际上，\(\int_0^\infty f(x) \, \mathrm{d}x\) 或 \(\int_0^\infty f(x) \sin^2 x \, \mathrm{d}x\) 发散的情形就是 (i)(ii) 的逆否命题。故结论得证。
\end{proof}

\begin{proposition}\label{proposition:反常积分敛散性运算性质}
设 \(f(x),g(x)\) 在任何闭区间上黎曼可积，其$f,g$在$x=a$处都有界.
\begin{enumerate}[(1)]
\item 若$\int_a^{\infty}{f\left( x \right) \mathrm{d}x}$绝对收敛,则$\int_a^{\infty}{f\left( x \right) \mathrm{d}x}$一定条件收敛.

\item 若$\int_a^{\infty}{f\left( x \right) \mathrm{d}x},
\int_a^{\infty}{g\left( x \right) \mathrm{d}x}$都绝对收敛,则$\int_a^{\infty}{\left[ f\left( x \right) \pm g\left( x \right) \right] \mathrm{d}x}$也绝对收敛.

\item 若$\int_a^{\infty}{f\left( x \right) \mathrm{d}x}$绝对收敛,$
\int_a^{\infty}{g\left( x \right) \mathrm{d}x}$条件收敛,则$\int_a^{\infty}{\left[ f\left( x \right) \pm g\left( x \right) \right] \mathrm{d}x}$条件收敛,但不绝对收敛.

\item 若$\int_a^{\infty}{f\left( x \right) \mathrm{d}x},
\int_a^{\infty}{g\left( x \right) \mathrm{d}x}$都条件收敛,则$\int_a^{\infty}{\left[ f\left( x \right) \pm g\left( x \right) \right] \mathrm{d}x}$可能条件收敛,也可能绝对收敛.
\end{enumerate}
\end{proposition}
\begin{proof}
\begin{enumerate}[(1)]
\item 由$f(x)\leqslant  |f(x)|$立得.

\item 由$|f(x)\pm g(x)|\leqslant  |f(x)|+|g(x)|$立得.

\item 由(1)可知$\int_a^{\infty}{f\left( x \right) \mathrm{d}x},
\int_a^{\infty}{g\left( x \right) \mathrm{d}x}$都条件收敛,从而$\int_a^{\infty}{\left[ f\left( x \right) \pm g\left( x \right) \right] \mathrm{d}x}$也条件收敛.若$\int_a^{\infty}{\left| f\left( x \right) \pm g\left( x \right) \right|\mathrm{d}x}$ $<\infty$,注意到$g\left( x \right) =\left[ f\left( x \right) +g\left( x \right) \right] -f\left( x \right)$,从而由(2)可知$\int_a^{\infty}{g\left( x \right) \mathrm{d}x}=\int_a^{\infty}{\left[ \left( f\left( x \right) +g\left( x \right) \right) -f\left( x \right) \right] \mathrm{d}x}$也绝对收敛,矛盾!

\item 
\end{enumerate}
\end{proof}

\begin{example}
判断如下积分的收敛性:
\begin{enumerate}
\item  \(\int_{0}^{\infty}\frac{1}{\sqrt[3]{x(x - 1)^2(x - 2)}}\mathrm{d}x\);

\item \(\int_{0}^{1}\frac{\sqrt[m]{\ln^2(1 - x)}}{\sqrt[n]{x}}\mathrm{d}x, m,n\in\mathbb{N}\);

\item \(\int_{2}^{\infty}(\sqrt{x + 1}-\sqrt{x})^p\ln\frac{x + 1}{x - 1}\mathrm{d}x\).
\end{enumerate}
\end{example}
\begin{proof}
\begin{enumerate}
\item 四个瑕点$x=0,1,2,\infty$,分别估阶讨论即得收敛.

\item 注意到
\begin{align*}
\frac{\sqrt[m]{\ln ^2(1 - x)}}{\sqrt[n]{x}}\sim\frac{x^{\frac{2}{m}}}{x^{\frac{1}{n}}}=\frac{1}{x^{\frac{1}{n}-\frac{2}{m}}},x\rightarrow 0^+.
\end{align*}
又\(\frac{1}{n}<1+\frac{2}{m},\forall m,n\in \mathbb{N}\),故\(\int_{0}^{\frac{1}{2}}\frac{\sqrt[m]{\ln ^2(1 - x)}}{\sqrt[n]{x}}\mathrm{d}x\left( \forall m,n\in \mathbb{N} \right)\)收敛。
注意到
\begin{align*}
\frac{\sqrt[m]{\ln ^2(1 - x)}}{\sqrt[n]{x}}\sim\ln ^{\frac{2}{m}}(1 - x),x\rightarrow 1^-.
\end{align*}
并且对\(\forall m\in \mathbb{N}\)，都有
\begin{align*}
\int_{\frac{1}{2}}^1{\ln ^{\frac{2}{m}}(1}-x)\mathrm{d}x&=\int_0^{\frac{1}{2}}{\ln ^{\frac{2}{m}}x\mathrm{d}x\xlongequal{x=e^t}\int_{-\infty}^{-\ln 2}{t^{\frac{2}{m}}e^t\mathrm{d}t}}
\\
&\xlongequal{t=-u}\int_{\ln 2}^{\infty}{u^{\frac{2}{m}}e^{-u}\mathrm{d}u}\leqslant \int_0^{\infty}{u^{\frac{2}{m}}e^{-u}\mathrm{d}u}
\\
&=\Gamma \left( 1+\frac{2}{m} \right) <+\infty .
\end{align*}
故\(\int_{\frac{1}{2}}^{1}\frac{\sqrt[m]{\ln ^2(1 - x)}}{\sqrt[n]{x}}\mathrm{d}x\left( \forall m,n\in \mathbb{N} \right)\)收敛。
综上，\(\int_{0}^{1}\frac{\sqrt[m]{\ln ^2(1 - x)}}{\sqrt[n]{x}}\mathrm{d}x\left( \forall m,n\in \mathbb{N} \right)\)收敛.

\item 由于\((\sqrt{x + 1} - \sqrt{x})^p\ln\frac{x + 1}{x - 1}\sim\frac{1}{x^{1+\frac{p}{2}}},x\rightarrow +\infty\)。故\(\int_{2}^{\infty}(\sqrt{x + 1} - \sqrt{x})^p\ln\frac{x + 1}{x - 1} \mathrm{d}x\)收敛当且仅当\(p > 0\)。 
\end{enumerate}
\end{proof}

\begin{example}
设\(p,q > 0\)，判断\(\int_{1}^{\infty}\frac{1}{x^p\ln^q x}\mathrm{d}x\)收敛性。 
\end{example}
\begin{note}
一个经验上的小结论. 在幂函数次数不为\(1\)时，趋于无穷或者趋于\(0\)时\(\ln\)可忽略.
\end{note}
\begin{proof}

先讨论\(\int_{1}^{2}\frac{1}{x^p\ln ^q x}\mathrm{d}x\)的收敛性. 由于
\begin{align*}
\frac{1}{x^p\ln ^q x}\sim\frac{1}{(x - 1)^q},x\rightarrow 1^+.
\end{align*}
因此\(\int_{1}^{2}\frac{1}{x^p\ln ^q x}\mathrm{d}x\)收敛当且仅当\(q < 1\).

再讨论\(\int_{2}^{\infty}\frac{1}{x^p\ln ^q x}\mathrm{d}x\)的收敛性.

\one 当\(p > 1\)时，我们有
\begin{align*}
\frac{\frac{1}{x^p\ln ^q x}}{\frac{1}{x^p}}=\frac{1}{\ln ^q x}\rightarrow 0,x\rightarrow +\infty.
\end{align*}
从而存在\(C > 0\)，当$x$充分大时,有
\begin{align*}
\frac{1}{x^p\ln ^q x}<\frac{C}{x^p}.
\end{align*}
而\(\int_{2}^{\infty}\frac{C}{x^p}\mathrm{d}x\)收敛，故\(\int_{2}^{\infty}\frac{1}{x^p\ln ^q x}\mathrm{d}x\)此时收敛.

\two 当\(0 < p < 1\)时，取\(\varepsilon > 0\)，使得\(p + \varepsilon < 1\)，从而
\begin{align*}
\frac{\frac{1}{x^p\ln ^q x}}{\frac{1}{x^{p + \varepsilon}}}=\frac{x^{\varepsilon}}{\ln ^q x}\rightarrow +\infty,x\rightarrow +\infty.
\end{align*}
于是存在\(M > 0\)，当$x$充分大时,有
\begin{align*}
\frac{1}{x^p\ln ^q x}>\frac{M}{x^{p + \varepsilon}}.
\end{align*}
而\(\int_{2}^{\infty}\frac{M}{x^{p + \varepsilon}}\mathrm{d}x\)发散，故\(\int_{2}^{\infty}\frac{1}{x^p\ln ^q x}\mathrm{d}x\)此时发散.

\three 当\(p = 1\)时，我们有
\begin{align*}
\int_{2}^{\infty}\frac{1}{x\ln ^q x}\mathrm{d}x=\int_{2}^{\infty}\frac{1}{\ln ^q x}\mathrm{d}\ln x=\int_{\ln 2}^{\infty}\frac{1}{t^q}\mathrm{d}t.
\end{align*}
于是此时\(\int_{2}^{\infty}\frac{1}{x\ln ^q x}\mathrm{d}x\)收敛当且仅当\(q > 1\).

综上所述，\(\int_{1}^{\infty}\frac{1}{x^p\ln ^q x}\mathrm{d}x\)收敛当且仅当\(p > 1,q < 1\). 
\end{proof}

\begin{example}
对\(a,b\in\mathbb{R}\)，判断\(\int_{0}^{\infty}\frac{\sin x^b}{x^a}\mathrm{d}x\)的收敛性和绝对收敛性. 
\end{example}
\begin{proof}

{\heiti 收敛性:}
\begin{enumerate}
\item 当\(b = 0\)时，此时\(\int_{0}^{\infty}\frac{\sin 1}{x^a}\mathrm{d}x\)必定发散.
\item 当\(b\neq 0\)时，我们有
\begin{align}\label{100.52}
\int_{0}^{+\infty}\frac{\sin x^b}{x^a}\mathrm{d}x&\xlongequal{y = x^b}\frac{1}{|b|}\int_{0}^{+\infty}\frac{\sin y}{y^{\frac{a}{b}}}y^{\frac{1}{b}-1}\mathrm{d}y=\frac{1}{|b|}\int_{0}^{+\infty}\frac{\sin y}{y^{\frac{a - 1}{b}+1}}\mathrm{d}y.
\end{align}
\begin{enumerate}
\item 先考虑\(\int_{0}^{1}\frac{\sin y}{y^{\frac{a - 1}{b}+1}}\mathrm{d}y\). 注意到
\begin{align*}
\frac{\sin y}{y^{\frac{a - 1}{b}+1}}\sim\frac{1}{y^{\frac{a - 1}{b}}},x\rightarrow 0^+.
\end{align*}
因此\(\int_{0}^{1}\frac{\sin y}{y^{\frac{a - 1}{b}+1}}\mathrm{d}y\)收敛当且仅当\(\frac{a - 1}{b}<1\).
\item 再考虑\(\int_{1}^{+\infty}\frac{\sin y}{y^{\frac{a - 1}{b}+1}}\mathrm{d}y\).
\begin{enumerate}
\item 当\(\frac{a - 1}{b}+1\leqslant 0\)时，我们有
\begin{align*}
&\left|\int_{\frac{\pi}{4}+2n\pi}^{\frac{\pi}{2}+2n\pi}\frac{\sin y}{y^{\frac{a - 1}{b}+1}}\mathrm{d}y\right| \geqslant \left(\frac{\pi}{4}+2n\pi\right)^{-\left(\frac{a - 1}{b}+1\right)}\left|\int_{\frac{\pi}{4}+2n\pi}^{\frac{\pi}{2}+2n\pi}\sin y\mathrm{d}y\right|\\
&=\left(\frac{\pi}{4}\right)^{-\left(\frac{a - 1}{b}+1\right)}\left|\int_{\frac{\pi}{4}}^{\frac{\pi}{2}}\sin y\mathrm{d}y\right|=\left(\frac{\pi}{4}\right)^{-\left(\frac{a - 1}{b}+1\right)}\left(1 - \frac{\sqrt{2}}{2}\right),\forall n\in\mathbb{N}.
\end{align*}
因此由Cauchy收敛准则可知，此时\(\int_{1}^{+\infty}\frac{\sin y}{y^{\frac{a - 1}{b}+1}}\mathrm{d}y\)发散.
\item 当\(\frac{a - 1}{b}+1>0\)时，我们有
\begin{align*}
\left| \int_0^x{\sin y\mathrm{d}y} \right|=\left| 1-\cos x \right|\leqslant 2\ (\forall x > 0),\quad \frac{1}{y^{\frac{a - 1}{b}+1}}\text{单调递减趋于}0\ (y\rightarrow +\infty).
\end{align*}
于是由Dirichlet判别法可知，此时\(\int_{1}^{+\infty}\frac{\sin y}{y^{\frac{a - 1}{b}+1}}\mathrm{d}y\)收敛.
\end{enumerate}
\end{enumerate}
\end{enumerate}
综上，\(\int_{0}^{\infty}\frac{\sin x^b}{x^a}\mathrm{d}x\)收敛当且仅当\(b\neq 0\)且\(-1<\frac{a - b}{b}<1\).

{\heiti 绝对收敛性:} 在\(-1 < \frac{a - 1}{b} < 1, b\neq 0\)情况下,先考虑$\int_0^1{\frac{\left| \sin y \right|}{y^{\frac{a-1}{b}+1}}\mathrm{d}y}$.我们有
\begin{align*}
\frac{\left| \sin y \right|}{y^{\frac{a-1}{b}+1}}\sim \frac{1}{y^{\frac{a-1}{b}}},x\rightarrow 0^+.
\end{align*}
又因为$\frac{a-1}{b}<1$,所以$\int_0^1{\frac{1}{y^{\frac{a-1}{b}}}\mathrm{d}y}$必收敛,因此$\int_0^1{\frac{\left| \sin y \right|}{y^{\frac{a-1}{b}+1}}\mathrm{d}y}$必绝对收敛.

再考虑$\int_{1}^{+\infty}{\frac{\left| \sin y \right|}{y^{\frac{a-1}{b}+1}}\mathrm{d}y}$.
当\(\frac{a - 1}{b}>0\)，注意到由\eqref{100.52}知道
\begin{align*}
\int_{1}^{+\infty}\frac{\sin x^b}{x^a}\mathrm{d}x=\frac{1}{|b|}\int_{1}^{\infty}\frac{|\sin y|}{y^{\frac{a - 1}{b}+1}}\mathrm{d}y \leqslant \frac{1}{|b|}\int_{1}^{\infty}\frac{1}{y^{\frac{a - 1}{b}+1}} < \infty,
\end{align*}
故此时$\int_{1}^{+\infty}{\frac{\left| \sin y \right|}{y^{\frac{a-1}{b}+1}}\mathrm{d}y}$绝对收敛.

当\(\frac{a - 1}{b}\leqslant 0\)，我们有
\begin{align*}
\frac{1}{|b|}\int_{1}^{\infty}\frac{|\sin y|}{y^{\frac{a - 1}{b}+1}}\mathrm{d}y &\geqslant \frac{1}{|b|}\int_{1}^{\infty}\frac{|\sin y|^2}{y^{\frac{a - 1}{b}+1}}\mathrm{d}y = \frac{1}{2|b|}\int_{1}^{\infty}\frac{1 - \cos(2y)}{y^{\frac{a - 1}{b}+1}}\mathrm{d}y \\
&= \frac{1}{2|b|}\int_{1}^{\infty}\frac{1}{y^{\frac{a - 1}{b}+1}}\mathrm{d}y - \frac{1}{2|b|}\int_{1}^{\infty}\frac{\cos(2y)}{y^{\frac{a - 1}{b}+1}}\mathrm{d}y
\end{align*}
因为$-1<\frac{a-1}{b}$,所以$\int_{1}^{\infty}\frac{1}{y^{\frac{a - 1}{b}+1}}\mathrm{d}y$发散,由Dirichlet判别法可知$\int_{1}^{\infty}\frac{\cos(2y)}{y^{\frac{a - 1}{b}+1}}\mathrm{d}y$收敛.故此时$\int_{1}^{+\infty}{\frac{\left| \sin y \right|}{y^{\frac{a-1}{b}+1}}\mathrm{d}y}$发散.

综上,这就证明了原积分在\(-1 < \frac{a - 1}{b} \leqslant 0, b\neq 0\)情况下条件收敛，\(0 < \frac{a - 1}{b} < 1, b\neq 0\)情况下绝对收敛. 
\end{proof}

\begin{example}
判断收敛性\(\int_{e^2}^{\infty}\frac{1}{\ln^{\ln x}x}\mathrm{d}x\)。 
\end{example}
\begin{note}
注意运用\(x^{\Box}=e^{\Box\ln x}=(e^{\Box})^{\ln x}\).
\end{note}
\begin{proof}
注意到
\begin{align*}
\int_{e^2}^{\infty}\frac{1}{\ln^{\ln x}x}\mathrm{d}x&=\int_{e^2}^{\infty}\frac{1}{e^{\ln x\cdot \ln\ln x}}\mathrm{d}x=\int_{e^2}^{\infty}\frac{1}{e^{\ln x\cdot \ln\ln x}}\mathrm{d}x\\
&=\int_{e^2}^{+\infty}\frac{1}{x^{\ln\ln x}}\mathrm{d}x=\int_{e^2}^{e^{e^2}}\frac{1}{x^{\ln\ln x}}\mathrm{d}x+\int_{e^{e^2}}^{\infty}\frac{1}{x^{\ln\ln x}}\mathrm{d}x\\
&\leqslant \int_{e^2}^{e^{e^2}}\frac{1}{x^{\ln\ln x}}\mathrm{d}x+\int_{e^{e^2}}^{\infty}\frac{1}{x^2}\mathrm{d}x<+\infty.
\end{align*}
故原积分收敛. 
\end{proof}

\begin{example}
判断收敛性和绝对收敛性:
\begin{enumerate}
\item \(\int_{1}^{\infty}\tan\left(\frac{\sin x}{x}\right)\mathrm{d}x\),

\item \(\int_{2}^{\infty}\frac{\sin x}{x^p(x^p + \sin x)}\mathrm{d}x,p > 0\). 
\end{enumerate}
\end{example}
\begin{note}
经验上,Taylor公式应该展开到余项里面的函数绝对收敛为止.
\end{note}
\begin{proof}
\begin{enumerate}
\item 由Taylor公式可知
\begin{align}
\tan\frac{\sin x}{x}=\frac{\sin x}{x}+O\left(\frac{\sin^3 x}{x^3}\right),x\rightarrow +\infty. \label{100.53}
\end{align}
由Dirichlet判别法可知\(\int_{1}^{\infty}\frac{\sin x}{x}\mathrm{d}x\)收敛,显然有\(\int_{1}^{\infty}\frac{\sin x}{x}\mathrm{d}x\)条件收敛. 注意到
\begin{align*}
\left|O\left(\frac{\sin^3 x}{x^3}\right)\right|\leqslant M\left|\frac{\sin x}{x}\right|^3\leqslant\frac{M}{x^3},x\rightarrow +\infty.
\end{align*}
故\(\int_{1}^{\infty}O\left(\frac{\sin^3 x}{x^3}\right)\mathrm{d}x\)绝对收敛. 因此由\eqref{100.53}式可得\(\int_{1}^{\infty}\tan\frac{\sin x}{x}\mathrm{d}x\)条件收敛.

\item 注意到
\begin{align*}
\int_{2}^{+\infty}\frac{\sin x}{x^p(x^p + \sin x)}\mathrm{d}x&=\int_{2}^{+\infty}\frac{\frac{\sin x}{x^p}}{x^p(1 + \frac{\sin x}{x^p})}\mathrm{d}x.
\end{align*}
取\(m\in \mathbb{N}\)，使\(m > \frac{1}{p}-1\)。由Taylor公式可知
\begin{align*}
\frac{t}{1 + t}=t - t^2+\cdots + (-1)^mt^{m - 1}+O(t^m),t\rightarrow 0^+.
\end{align*}
从而
\begin{align}
\frac{\frac{\sin x}{x^p}}{x^p(1 + \frac{\sin x}{x^p})}=\frac{\sin x}{x^{2p}}-\frac{\sin^2 x}{x^{3p}}+\cdots + (-1)^m\frac{\sin^{m - 1} x}{x^{mp}}+O\left(\frac{\sin^m x}{x^{(m + 1)p}}\right),x\rightarrow +\infty. \label{100.58}
\end{align}
注意到
\begin{align}
\frac{\sin^2 x}{x^{3p}}=\frac{1}{2x^{3p}}-\frac{\cos 2x}{x^{3p}}, \label{100.57}
\end{align}
\begin{enumerate}
\item[(i)] 当\(p\leqslant \frac{1}{3}\)时，有\(\int_{2}^{+\infty}\frac{1}{2x^{3p}}\mathrm{d}x\)发散，从而此时\(\int_{2}^{+\infty}\frac{\sin x}{x^p(x^p + \sin x)}\mathrm{d}x\)发散.
\item[(ii)] 当\(p > \frac{1}{3}\)时，有\(\int_{2}^{+\infty}\frac{1}{2x^{3p}}\mathrm{d}x\)收敛，并且由Dirichlet判别法可知\(\int_{2}^{+\infty}\frac{\sin x}{x^{2p}}\mathrm{d}x\)，\(\int_{2}^{+\infty}\frac{\cos 2x}{x^{3p}}\mathrm{d}x\)收敛。从而由\eqref{100.57}式可知，此时\(\int_{2}^{+\infty}\frac{\sin^2 x}{x^{3p}}\mathrm{d}x\)收敛.
又因为对\(\forall k\geqslant 2\)，都有
\begin{align*}
\left|\frac{\sin^k x}{x^{(k + 1)p}}\right|\leqslant\frac{1}{x^{(k + 1)p}}\leqslant\frac{1}{x^{3p}},\forall x > 2.
\end{align*}
而\(\int_{2}^{+\infty}\frac{1}{x^{3p}}\mathrm{d}x\)收敛，所以此时\(\int_{2}^{+\infty}\frac{\sin^k x}{x^{(k + 1)p}}\mathrm{d}x\)都绝对收敛。故由\eqref{100.58}式可知，此时原积分收敛.
\end{enumerate}
综上，原积分在\(p\leqslant \frac{1}{3}\)时发散，\(p > \frac{1}{3}\)收敛.
再讨论绝对收敛性.
\begin{enumerate}
\item 当\(p > \frac{1}{2}\)时，由$M$判别法易知\(\int_{2}^{+\infty}\frac{\sin^k x}{x^{(k + 1)p}}\mathrm{d}x(1\leqslant  k\leqslant  m)\)绝对收敛。再由\eqref{100.58}式可知，此时原积分绝对收敛.
\item 当\(\frac{1}{3}<p\leqslant \frac{1}{2}\)时，我们有
\begin{align*}
\left|\frac{\sin x}{x^{2p}}\right|\geqslant\frac{\sin^2 x}{x^{2p}}=\frac{1}{2x^{2p}}-\frac{\cos 2x}{2x^{2p}}.
\end{align*}
显然\(\int_{2}^{+\infty}\frac{1}{2x^{2p}}\mathrm{d}x\)发散，故此时\(\int_{2}^{+\infty}\frac{\sin x}{x^{2p}}\mathrm{d}x\)条件收敛。注意到对\(\forall k\geqslant 2\)，都有
\begin{align*}
\left|\frac{\sin^k x}{x^{(k + 1)p}}\right|\leqslant\frac{1}{x^{(k + 1)p}}\leqslant\frac{1}{x^{3p}},\forall x > 2.
\end{align*}
而显然此时\(\int_{2}^{+\infty}\frac{1}{x^{3p}}\mathrm{d}x\)收敛，故\(\int_{2}^{+\infty}\frac{\sin^k x}{x^{(k + 1)p}}\mathrm{d}x(2\leqslant  k\leqslant  m)\)绝对收敛。因此再由\eqref{100.58}式及\nrefpro{proposition:反常积分敛散性运算性质}{(3)}可知原积分此时条件收敛.
\end{enumerate} 
\end{enumerate}
\end{proof}

\begin{example}
设\(f(x)\)在\(\mathbb{R}\)上非负连续，对任意正整数\(k\)有$\int_{-\infty}^{+\infty} e^{-\frac{|x|}{k}}f(x)\mathrm{d}x \leqslant  1,$
证明：$\int_{-\infty}^{+\infty}f(x)\mathrm{d}x \leqslant  1.$
\end{example}
\begin{remark}
实际上,由实变函数相关结论可直接得到
\begin{align*}
\lim_{k\rightarrow \infty}\int_{\mathbb{R}}e^{-\frac{| x |}{k}}f( x ) \mathrm{d}x=\int_{\mathbb{R}}\left[ \lim_{k\rightarrow \infty}e^{-\frac{| x |}{k}}f( x ) \right] \mathrm{d}x=\int_{\mathbb{R}}f( x ) \mathrm{d}x.
\end{align*}
\end{remark}
\begin{proof}
由条件可得，对$\forall A>0$，我们有
\begin{align*}
1\geqslant \int_{-A}^A{e^{-\frac{|x|}{k}}f(x)\mathrm{d}x}\geqslant e^{-\frac{1}{k}}\int_{-A}^A{f(x)\mathrm{d}x}\Longrightarrow \int_{-A}^A{f(x)\mathrm{d}x}\leqslant e^{-\frac{1}{k}},\forall k\in \mathbb{N} .
\end{align*}
令$k\rightarrow \infty$，则$\int_{-A}^Af( x ) \mathrm{d}x\leqslant 1,\forall A>0$。于是再令$A\rightarrow +\infty$，可得$\int_{\mathbb{R}}f( x ) \mathrm{d}x\leqslant 1.$

实际上再由单调有界可知$\int_{\mathbb{R}}f( x ) \mathrm{d}x$收敛。
\end{proof}

\begin{example}
对实数\(a\)，讨论$\int_{0}^{\infty}\frac{x}{\cos^{2}x + x^{a}\sin^{2}x}\mathrm{d}x$
的敛散性. 
\end{example}
\begin{proof}

{\color{blue}证法一:}
先讨论$\int_0^1\frac{x}{\cos^2x+x^a\sin^2x}\mathrm{d}x$的敛散性. 注意到
\begin{align*}
\int_0^1\frac{x}{\cos^2x+x^a\sin^2x}\mathrm{d}x\leqslant \int_0^1\frac{1}{\cos^2x}\mathrm{d}x=\tan x \Big |_{0}^{1}=\tan 1<\infty,\quad \forall a\in \mathbb{R}.
\end{align*}
故$\forall a\in \mathbb{R}$，都有$\int_0^1\frac{x}{\cos^2x+x^a\sin^2x}\mathrm{d}x$收敛.
再讨论$\int_1^{\infty}\frac{x}{\cos^2x+x^a\sin^2x}\mathrm{d}x$的敛散性.
\begin{enumerate}[(i)]
\item 当$a\leqslant 2$时，
\begin{align*}
\int_1^{\infty}\frac{x}{\cos^2x+x^a\sin^2x}\mathrm{d}x\geqslant \int_1^{\infty}\frac{x}{\cos^2x+x^2\sin^2x}\mathrm{d}x\geqslant \int_1^{\infty}\frac{x}{x^2+1}\mathrm{d}x=\frac{1}{2}\int_1^{\infty}\frac{1}{x^2+1}\mathrm{d}(x^2+1)=+\infty.
\end{align*}
\item 当$a>2$时，我们有（等价关系直观上是显然的，可由拟合法或放缩严谨证明）
\begin{align}
\int_{n\pi}^{(n+1)\pi}\frac{x}{\cos^2x+x^a\sin^2x}\mathrm{d}x=\int_0^{\pi}\frac{x+n\pi}{\cos^2x+(x+n\pi)^a\sin^2x}\mathrm{d}x\sim n\pi \int_0^{\pi}\frac{1}{\cos^2x+(n\pi)^a\sin^2x}\mathrm{d}x,\quad n\rightarrow \infty.\label{eq:100.34}
\end{align}
注意到对$\forall \lambda >0$,我们都有
\begin{align*}
\int_0^{\pi}\frac{1}{\cos^2x+\lambda \sin^2x}\mathrm{d}x=2\int_0^{\frac{\pi}{2}}\frac{1}{\cos^2x+\lambda \sin^2x}\mathrm{d}x=2\int_0^{\frac{\pi}{2}}\frac{1}{1+\lambda \tan^2x}\cdot \frac{1}{\cos^2x}\mathrm{d}x=2\int_0^{\infty}\frac{1}{1+\lambda t^2}\mathrm{d}t=\frac{\pi}{\sqrt{\lambda}}.
\end{align*}
故再结合\eqref{eq:100.34}式可知
\begin{align*}
\int_{n\pi}^{(n+1)\pi}\frac{x}{\cos^2x+x^a\sin^2x}\mathrm{d}x\sim n\pi \int_0^{\pi}\frac{1}{\cos^2x+(n\pi)^a\sin^2x}\mathrm{d}x\sim n\pi \frac{\pi}{(n\pi)^{\frac{a}{2}}}\sim \frac{1}{n^{\frac{a}{2}-1}},\quad n\rightarrow \infty.
\end{align*}
由于被积函数非负,因此再结合\refpro{proposition:反常积分收敛的基本结论2}可知
\begin{align*}
\int_1^{\infty}\frac{x}{\cos^2x+x^a\sin^2x}\mathrm{d}x\sim \int_{\pi}^{\infty}\frac{x}{\cos^2x+x^a\sin^2x}\mathrm{d}x=\sum_{n=1}^{\infty}\int_{n\pi}^{(n+1)\pi}\frac{x}{\cos^2x+x^a\sin^2x}\mathrm{d}x\sim \sum_{n=1}^{\infty}\frac{1}{n^{\frac{a}{2}-1}},\quad n\rightarrow \infty.
\end{align*}
从而当$\frac{a}{2}-1\leqslant 1$时，即$2<a\leqslant 4$，$\int_1^{\infty}\frac{x}{\cos^2x+x^a\sin^2x}\mathrm{d}x$发散；当$\frac{a}{2}-1>1$，即$a>4$时，$\int_1^{\infty}\frac{x}{\cos^2x+x^a\sin^2x}\mathrm{d}x$收敛.
\end{enumerate}

综上，当$a>4$时，$\int_0^{\infty}\frac{x}{\cos^2x+x^a\sin^2x}\mathrm{d}x$收敛；当$a\leqslant 4$时，$\int_0^{\infty}\frac{x}{\cos^2x+x^a\sin^2x}\mathrm{d}x$发散.

{\color{blue}证法二:}由于被积函数非负,因此由\refpro{proposition:反常积分收敛的基本结论2}可知
\begin{align*}
\int_0^{+\infty}{\frac{x}{\cos ^2x+x^a\sin ^2x}\mathrm{d}x}&=\underset{n\rightarrow \infty}{\lim}\int_0^{n\pi}{\frac{x}{\cos ^2x+x^a\sin ^2x}\mathrm{d}x}
=\sum_{n=1}^{\infty}{\int_{\left( n-1 \right) \pi}^{n\pi}{\frac{x}{\cos ^2x+x^a\sin ^2x}\mathrm{d}x}}
\\
&=\sum_{n=1}^{\infty}{\int_0^{\pi}{\frac{\left( n-1 \right) \pi +y}{\cos ^2y+\left[ \left( n-1 \right) \pi +y \right] ^a\sin ^2y}\mathrm{d}y}}.
\end{align*}
一方面，我们有
\begin{align*}
\sum_{n=1}^{\infty}{\int_0^{\pi}{\frac{\left( n-1 \right) \pi +y}{\cos ^2y+\left[ \left( n-1 \right) \pi +y \right] ^a\sin ^2y}\mathrm{d}y}}&\leqslant \sum_{n=1}^{\infty}{\int_0^{\pi}{\frac{n\pi}{\cos ^2y+\left[ \left( n-1 \right) \pi \right] ^a\sin ^2y}\mathrm{d}y}} =\sum_{n=1}^{\infty}{\int_0^{\frac{\pi}{2}}{\frac{2n\pi}{\cos ^2y+\left[ \left( n-1 \right) \pi \right] ^a\sin ^2y}\mathrm{d}y}}\\
&=\sum_{n=1}^{\infty}{\int_0^{\frac{\pi}{2}}{\frac{\frac{2n\pi}{\cos ^2y}}{1+\left[ \left( n-1 \right) \pi \right] ^a\tan ^2y}\mathrm{d}y}}
\xlongequal{t=\tan y}2\pi \sum_{n=1}^{\infty}{n\int_0^{+\infty}{\frac{\mathrm{d}t}{1+\left[ \left( n-1 \right) \pi \right] ^at^2}}}\\
&=\pi ^2\sum_{n=1}^{\infty}{\frac{n}{\sqrt{\left[ \left( n-1 \right) \pi \right] ^a}}}\sim \sum_{n=1}^{\infty}{\frac{C_1}{n^{\frac{a}{2}-1}}},n\rightarrow \infty .
\end{align*}
故当$a<4$时，有$\frac{a}{2}-1<1$，此时原积分收敛。另一方面，我们有
\begin{align*}
\sum_{n=1}^{\infty}{\int_0^{\pi}{\frac{\left( n-1 \right) \pi +y}{\cos ^2y+\left[ \left( n-1 \right) \pi +y \right] ^a\sin ^2y}\mathrm{d}y}}&\geqslant \sum_{n=1}^{\infty}{\int_0^{\pi}{\frac{\left( n-1 \right) \pi}{\cos ^2y+\left( n\pi \right) ^a\sin ^2y}\mathrm{d}y}}=\sum_{n=1}^{\infty}{\int_0^{\frac{\pi}{2}}{\frac{2\left( n-1 \right) \pi}{\cos ^2y+\left( n\pi \right) ^a\sin ^2y}\mathrm{d}y}}\\
&=\sum_{n=1}^{\infty}{\int_0^{\frac{\pi}{2}}{\frac{\frac{2\left( n-1 \right) \pi}{\cos ^2y}}{1+\left( n\pi \right) ^a\tan ^2y}\mathrm{d}y}}\xlongequal{t=\tan y}2\pi \sum_{n=1}^{\infty}{\left( n-1 \right) \int_0^{+\infty}{\frac{\mathrm{d}t}{1+\left( n\pi \right) ^at^2}}}\\
&=\pi ^2\sum_{n=1}^{\infty}{\frac{n-1}{\sqrt{\left( n\pi \right) ^a}}}\sim \sum_{n=1}^{\infty}{\frac{C_2}{n^{\frac{a}{2}-1}}},n\rightarrow \infty .
\end{align*}
故当$a\geqslant 4$时，有$\frac{a}{2}-1\geqslant 1$，此时原积分发散。
\end{proof}

\begin{example}
对$x > 0$, 判断积分$\int_0^{\infty} \frac{[t] - t + a}{t + x} \mathrm{d}t$收敛性.
\end{example}
\begin{proof}
注意到
\begin{align}
\lim_{n\to\infty}\int_1^n\frac{[t] - t + a}{t + x}\mathrm{d}t &= \lim_{n\to\infty}\sum_{k=1}^{n-1}\int_k^{k+1}\frac{k - t + a}{t + x}\mathrm{d}t = \sum_{k=1}^{\infty}\int_k^{k+1}\frac{k - t + a}{t + x}\mathrm{d}t\nonumber\\
&= \sum_{k=1}^{\infty}\left[(a + k + x)\ln\left(1 + \frac{1}{k + x}\right) - 1\right] = \sum_{k=1}^{\infty}\left[\frac{a - \frac{1}{2}}{k} + O\left(\frac{1}{k^2}\right)\right].\label{100.115}
\end{align}
故当$a \ne \frac{1}{2}$时，有$\sum_{k=1}^{\infty}\frac{a - \frac{1}{2}}{k}$发散，从而结合\eqref{100.115}式可知，此时$\lim_{n\to\infty}\int_1^n\frac{[t] - t + \frac{1}{2}}{t + x}\mathrm{d}t$不存在。因此由\hyperref[proposition:子列极限命题]{子列极限命题(a)}可知，此时$\int_0^{\infty}\frac{[t] - t + a}{t + x}\mathrm{d}t$发散。

当$a = \frac{1}{2}$时，由\eqref{100.115}式知
\begin{align}
\lim_{n\to\infty}\int_1^n\frac{[t] - t + \frac{1}{2}}{t + x}\mathrm{d}t = \sum_{k=1}^{\infty}O\left(\frac{1}{k^2}\right) \leqslant M\sum_{k=1}^{\infty}\frac{1}{k^2} < \infty.\label{100.114}
\end{align}
对$\forall y > 0$，存在唯一$n \in \mathbb{N}$，使得$n \leqslant y < n + 1$。于是
\begin{align}
\int_0^y\frac{[t] - t + \frac{1}{2}}{t + x}\mathrm{d}t = \int_0^n\frac{[t] - t + \frac{1}{2}}{t + x}\mathrm{d}t + \int_n^y\frac{[t] - t + \frac{1}{2}}{t + x}\mathrm{d}t.\label{100.113}
\end{align}
注意到
\begin{align*}
\left|\int_n^y\frac{[t] - t + \frac{1}{2}}{t + x}\mathrm{d}t\right| \leqslant \int_n^y\frac{1 + \frac{1}{2}}{t + x}\mathrm{d}t \leqslant \frac{3}{2}\int_n^{n+1}\frac{1}{t + x}\mathrm{d}t = \frac{3}{2}\ln\frac{n + 1 + x}{n + x}.
\end{align*}
当$y \to +\infty$时，有$n \to +\infty$，故$\int_n^y\frac{[t] - t + \frac{1}{2}}{t + x}\mathrm{d}t \to 0$，$y \to +\infty$。再结合\eqref{100.114}\eqref{100.113}式可知
\begin{align*}
\int_0^{\infty}\frac{[t] - t + \frac{1}{2}}{t + x}\mathrm{d}t = \lim_{n\to\infty}\int_0^n\frac{[t] - t + \frac{1}{2}}{t + x}\mathrm{d}t < \infty.
\end{align*}
故当$a = \frac{1}{2}$时，$\int_0^{\infty}\frac{[t] - t + \frac{1}{2}}{t + x}\mathrm{d}t$收敛。
\end{proof}

\begin{example}
对正整数\(n\)，讨论$\int_{0}^{+\infty}x^{n}e^{-x^{12}\sin^{2}x}\mathrm{d}x$
的敛散性. 
\end{example}
\begin{proof}
注意到
\begin{align}\label{equation:100.35}
\int_{k\pi}^{(k+1)\pi}x^ne^{-x^{12}\sin^2x}\mathrm{d}x=\int_0^{\pi}(x+k\pi)^ne^{-(x+k\pi)^{12}\sin^2x}\mathrm{d}x\sim (k\pi)^n\int_0^{\pi}e^{-(x+k\pi)^{12}\sin^2x}\mathrm{d}x,\quad k\rightarrow \infty.
\end{align}
又注意到
\begin{align*}
\int_0^{\pi}e^{-\lambda \sin^2x}\mathrm{d}x=2\int_0^{\frac{\pi}{2}}e^{-\lambda \sin^2x}\mathrm{d}x\geqslant 2\int_0^{\frac{\pi}{2}}e^{-\lambda x^2}\mathrm{d}x=\frac{2}{\sqrt{\lambda}}\int_0^{\frac{\pi}{2}\sqrt{\lambda}}e^{-x^2}\mathrm{d}x\sim \sqrt{\frac{\pi}{\lambda}},\quad \lambda \rightarrow +\infty,
\end{align*}
\begin{align*}
\int_0^{\pi}e^{-\lambda \sin^2x}\mathrm{d}x=2\int_0^{\frac{\pi}{2}}e^{-\lambda \sin^2x}\mathrm{d}x\leqslant 2\int_0^{\frac{\pi}{2}}e^{-\lambda \frac{4}{\pi^2}x^2}\mathrm{d}x=\frac{\pi}{\sqrt{\lambda}}\int_0^{\sqrt{\lambda}}e^{-x^2}\mathrm{d}x\sim \frac{\pi \sqrt{\pi}}{2\sqrt{\lambda}},\quad \lambda \rightarrow +\infty.
\end{align*}
故$\int_0^{\pi}e^{-\lambda \sin^2x}\mathrm{d}x\sim \frac{C}{\sqrt{\lambda}}$，$\lambda \rightarrow +\infty$，其中$C$为某一常数. 因此
\begin{align*}
\int_0^{\pi}e^{-(k\pi)^{12}\sin^2x}\mathrm{d}x\sim \frac{C}{(k\pi)^6},\quad k\rightarrow +\infty,
\end{align*}
\begin{align*}
\int_0^{\pi}e^{-[(k+1)\pi]^{12}\sin^2x}\mathrm{d}x\sim \frac{C}{[(k+1)\pi]^6},\quad k\rightarrow +\infty.
\end{align*}
又因为
\begin{align*}
\int_0^{\pi}e^{-[(k+1)\pi]^{12}\sin^2x}\mathrm{d}x\leqslant \int_0^{\pi}e^{-(x+k\pi)^{12}\sin^2x}\mathrm{d}x\leqslant \int_0^{\pi}e^{-(k\pi)^{12}\sin^2x}\mathrm{d}x,
\end{align*}
所以$\int_0^{\pi}e^{-(x+k\pi)^{12}\sin^2x}\mathrm{d}x\sim \frac{C_1}{k^6}$，$k\rightarrow +\infty$，其中$C_1$为某一常数. 于是结合\eqref{equation:100.35}式可知
\begin{align*}
\int_{k\pi}^{(k+1)\pi}x^ne^{-x^{12}\sin^2x}\mathrm{d}x=\int_0^{\pi}(x+k\pi)^ne^{-(x+k\pi)^{12}\sin^2x}\mathrm{d}x\sim (k\pi)^n\int_0^{\pi}e^{-(x+k\pi)^{12}\sin^2x}\mathrm{d}x\sim C_2k^{n-6},\quad k\rightarrow \infty.
\end{align*}
其中$C_2$为某一常数. 因此再结合\refpro{proposition:反常积分收敛的基本结论2}可得
\begin{align*}
\int_{\pi}^{\infty}{x^ne^{-x^{12}\sin ^2x}\mathrm{d}x}=\underset{k\rightarrow \infty}{\lim}\int_0^{k\pi}{x^ne^{-x^{12}\sin ^2x}\mathrm{d}x}=\sum_{k=1}^{\infty}{\int_{k\pi}^{(k+1)\pi}{x^ne^{-x^{12}\sin ^2x}\mathrm{d}x}}\sim \sum_{k=1}^{\infty}{C_2k^{n-6},\quad k\rightarrow \infty}.
\end{align*}
故当$n<5$时，$\int_{\pi}^{\infty}x^ne^{-x^{12}\sin^2x}\mathrm{d}x$收敛；当$n\geqslant 5$时，$\int_{\pi}^{\infty}x^ne^{-x^{12}\sin^2x}\mathrm{d}x$发散. 又因为
\begin{align*}
\int_0^{\pi}x^ne^{-x^{12}\sin^2x}\mathrm{d}x\leqslant \pi^n,
\end{align*}
所以$\int_0^{\pi}x^ne^{-x^{12}\sin^2x}\mathrm{d}x$对$\forall n\in \mathbb{N}$都收敛. 从而由
\begin{align*}
\int_0^{\infty}x^ne^{-x^{12}\sin^2x}\mathrm{d}x=\int_0^{\pi}x^ne^{-x^{12}\sin^2x}\mathrm{d}x+\int_{\pi}^{\infty}x^ne^{-x^{12}\sin^2x}\mathrm{d}x,
\end{align*}
可知当$n<5$时，$\int_0^{\infty}x^ne^{-x^{12}\sin^2x}\mathrm{d}x$收敛；当$n\geqslant 5$时，$\int_0^{\infty}x^ne^{-x^{12}\sin^2x}\mathrm{d}x$发散.
\end{proof}

\begin{example}
设\(p,q\)为正整数，求反常积分$I(p,q)=\int_{0}^{+\infty}\frac{\cos^{p}x - \cos^{q}x}{x}\mathrm{d}x$
收敛的充要条件.
\end{example}
\begin{proof}
因为当$p=q$时,积分显然收敛,所以只需考虑$p\ne q$的情形. 由$I(q,p) =-I(p,q)$可知,可以不妨设$p>q$,否则用$I(q,p) =-I(p,q)$代替$I(p,q)$即可

先讨论$\int_0^1{\frac{\cos ^px-\cos ^qx}{x}\mathrm{d}x}$的敛散性.由Taylor定理可知,对$\forall \varepsilon \in (0,1)$,存在$\delta >0$ ,使得  
\[
-\frac{x^2}{2}-\varepsilon x^2\leqslant \cos x\leqslant 1-\frac{x^2}{2}+\varepsilon x^2 ,\quad \forall x\in [0,\delta].
\]  
于是  
\begin{align*}
\int_0^1{\frac{\cos ^px-\cos ^qx}{x}\mathrm{d}x}&=\int_0^{\delta}{\frac{\cos ^px-\cos ^qx}{x}\mathrm{d}x}+\int_{\delta}^1{\frac{\cos ^px-\cos ^qx}{x}\mathrm{d}x}\\
&\leqslant \int_0^{\delta}{\frac{(1-\frac{x^2}{2}+\varepsilon x^2)^p-(1-\frac{x^2}{2}-\varepsilon x^2)^q}{x}\mathrm{d}x}+\frac{2}{\delta}(1-\delta)\\
&\leqslant \int_0^{\delta}{\frac{\frac{q-p+(p-q)\varepsilon}{2}x^2+(p+q)\mathrm{C}_{p}^{2}x^4}{x}\mathrm{d}x}+\frac{2}{\delta}(1-\delta)\\
&=\frac{q-p+(p-q)\varepsilon}{4}\delta +\frac{(p+q)\mathrm{C}_{p}^{2}}{4}\delta +\frac{2}{\delta}(1-\delta).
\end{align*}  
令$\varepsilon \rightarrow 0^+$,得$\int_0^1{\frac{\cos ^px-\cos ^qx}{x}\mathrm{d}x}\leqslant \frac{q-p}{4}\delta +\frac{(p+q)\mathrm{C}_{p}^{2}}{4}\delta +\frac{2}{\delta}(1-\delta)$.故对$\forall p>q$且$p,q\in \mathbb{N}$,都有$\int_0^1{\frac{\cos ^px-\cos ^qx}{x}\mathrm{d}x}$收敛.

再讨论$\int_1^{\infty}{\frac{\cos ^px-\cos ^qx}{x}\mathrm{d}x}$的敛散性.

(i) 当$p,q$都是奇数时,由\refthe{theorem:sinnx和cosnx以及sin^nx和cos^nx的展开式}可知  
$$
\cos ^px=\sum_{k=1}^p{p_k\cos kx} ,\quad \text{其中}p_k\in \mathbb{R} , k=1,2,\cdots,p.
$$  
$$
\cos ^qx=\sum_{k=1}^q{q_k\cos kx} ,\quad \text{其中}q_k\in \mathbb{R} , k=1,2,\cdots,q.
$$  
从而此时  
\begin{align*}
\int_1^{\infty}{\frac{\cos ^px-\cos ^qx}{x}\mathrm{d}x}&=\int_1^{\infty}{\frac{\sum\limits_{k=1}^p{p_k\cos kx}-\sum\limits_{k=1}^q{q_k\cos kx}}{x}\mathrm{d}x}\\
&=\sum\limits_{k=1}^q{(p_k-q_k)\int_1^{\infty}{\frac{\cos kx}{x}\mathrm{d}x}}+\sum\limits_{k=q+1}^p{p_k\int_1^{\infty}{\frac{\cos kx}{x}\mathrm{d}x}} .\label{100.36}
\end{align*}  
注意到对$\forall k\in \mathbb{N}$ 都有  
$$
\int_1^x{\cos kt\mathrm{d}t}=\frac{\sin kx-\sin k}{k}<2 ,\quad \forall x>1.
$$  
并且$\frac{1}{x}$在$[1,+\infty)$上单调递减趋于0,故由Dirichlet判别法可知,$\int_1^{\infty}{\frac{\cos kx}{x}\mathrm{d}x}(k\in \mathbb{N})$都收敛.因此再结合\eqref{100.36}式可知,$\int_1^{\infty}{\frac{\cos ^px-\cos ^qx}{x}\mathrm{d}x}$收敛 . 

(ii) 当$p,q$中至少有一个是偶数时,不妨设$p$是偶数 $q$不是偶数 ,则由\refthe{theorem:sinnx和cosnx以及sin^nx和cos^nx的展开式}可知  
\[
\cos ^px=\frac{1}{2^{p-1}}\sum_{k=0}^{\frac{p}{2}-1}{\mathrm{C}_{p}^{k}\cos 2\left( \frac{p}{2}-k \right) x}+\frac{\mathrm{C}_{p}^{\frac{p}{2}}}{2^p}.
\]  
\[
\cos ^qx=\sum_{k=1}^q{q_k\cos kx} \quad \text{其中}q_k\in \mathbb{R} \quad k=1,2,\cdots,q.
\] 
于是
\begin{align*}
\int_1^{\infty}{\frac{\cos ^px-\cos ^qx}{x}\mathrm{d}x}&=\int_1^{\infty}{\frac{\frac{1}{2^{p-1}}\sum\limits_{k=0}^{\frac{p}{2}-1}{\mathrm{C}_{p}^{k}\cos 2\left( \frac{p}{2}-k \right) x}-\sum\limits_{k=1}^q{q_k\cos kx}+\frac{\mathrm{C}_{p}^{\frac{p}{2}}}{2^p}}{x}\mathrm{d}x}\\
&=\int_1^{\infty}{\frac{\frac{1}{2^{p-1}}\sum\limits_{k=0}^{\frac{p}{2}-1}{\mathrm{C}_{p}^{k}\cos 2\left( \frac{p}{2}-k \right) x}-\sum\limits_{k=1}^q{q_k\cos kx}}{x}\mathrm{d}x}+\frac{\mathrm{C}_{p}^{\frac{p}{2}}}{2^p}\int_1^{\infty}{\frac{1}{x}\mathrm{d}x}.
\end{align*}  
由于$\int_1^{\infty}{\frac{1}{x}\mathrm{d}x}$发散, 故此时$\int_1^{\infty}{\frac{\cos ^px-\cos ^qx}{x}\mathrm{d}x}$也发散  .

综上,由  
\[
\int_0^{\infty}{\frac{\cos ^px-\cos ^qx}{x}\mathrm{d}x}=\int_0^1{\frac{\cos ^px-\cos ^qx}{x}\mathrm{d}x}+\int_1^{\infty}{\frac{\cos ^px-\cos ^qx}{x}\mathrm{d}x}.
\]
可知当$p=q$或$p,q$均为奇数时, $\int_0^{\infty}{\frac{\cos ^px-\cos ^qx}{x}\mathrm{d}x}$收敛 ,其余情形均发散.
\end{proof}

\begin{example}
对实数$p\ne 0$,讨论$I = \int_0^1{\frac{\cos(\frac{1}{1 - x})}{\sqrt[p]{1 - x^2}}\mathrm{d}x}$的敛散性.
\end{example}
\begin{proof}
对$I$进行积分换元可得
\begin{align}
I=\int_0^1{\frac{\cos\left(\frac{1}{1-x}\right)}{\sqrt[p]{1-x^2}}\mathrm{d}x}\xlongequal{u=\frac{1}{1-x}}\int_1^{\infty}{\frac{\cos u}{\left( 1-\left( 1-\frac{1}{u} \right) ^2 \right) ^{\frac{1}{p}}}\cdot \frac{1}{u^2}\mathrm{d}u} \notag \\
=\int_1^{\infty}{\frac{\cos u}{\left( \frac{2}{u}-\frac{1}{u^2} \right) ^{\frac{1}{p}}u^2}\mathrm{d}u}=\int_1^{\infty}{\frac{\cos u}{\left( 2-\frac{1}{u} \right) ^{\frac{1}{p}}u^{2-\frac{1}{p}}}\mathrm{d}u}.\label{eq:integral-transform111}
\end{align}

(i) 当$p>\frac{1}{2}$时, 令$f\left( u \right) =\left[ \left( 2-\frac{1}{u} \right) ^{\frac{1}{p}}u^{2-\frac{1}{p}} \right] ^p=\left( 2-\frac{1}{u} \right) u^{2p-1}$, 则显然有$\underset{u\rightarrow +\infty}{\lim}f\left( u \right) =+\infty$且$f\left( u \right)$递增. 于是$\frac{1}{\left( \frac{2}{u}-\frac{1}{u^2} \right) ^{\frac{1}{p}}u^2}=\frac{1}{\sqrt[p]{f\left( u \right)}}$在$\left[ 1,+\infty \right)$上单调递减趋于$0$. 又显然有$\int_1^A{\cos x\mathrm{d}x}$关于$A$有界, 所以结合\eqref{eq:integral-transform111}式, 再由Dirichlet判别法可知$I$收敛.

(ii) 当$p\in \left[ 0,\frac{1}{2} \right]$时, 若$p=\frac{1}{2}$, 则$\underset{u\rightarrow +\infty}{\lim}\frac{1}{\left( 2-\frac{1}{u} \right) ^{\frac{1}{p}}u^{2-\frac{1}{p}}}=2$; 若$p\in \left[ 0,\frac{1}{2} \right)$, 则$\underset{u\rightarrow +\infty}{\lim}\frac{1}{\left( 2-\frac{1}{u} \right) ^{\frac{1}{p}}u^{2-\frac{1}{p}}}=+\infty$. 因此对$\forall p\in \left[ 0,\frac{1}{2} \right]$, 都存在$K>0$, 使得
$$\frac{1}{\left( 2-\frac{1}{u} \right) ^{\frac{1}{p}}u^{2-\frac{1}{p}}}\geqslant 1,\forall u>K.$$
于是对$\forall k\in \mathbb{N} \cap \left( K,+\infty \right)$, 都有
$$\left| \int_{\frac{k\pi}{2}}^{\frac{\left( k+1 \right) \pi}{2}}{\frac{\cos u}{\left( 2-\frac{1}{u} \right) ^{\frac{1}{p}}u^{2-\frac{1}{p}}}\mathrm{d}u} \right|\geqslant \left| \int_{\frac{k\pi}{2}}^{\frac{\left( k+1 \right) \pi}{2}}{\cos u\mathrm{d}u} \right|=1.$$
故由Cauchy收敛准则可知,$I=\int_1^{\infty}{\frac{\cos u}{\left( 2-\frac{1}{u} \right) ^{\frac{1}{p}}u^{2-\frac{1}{p}}}\mathrm{d}u}$发散.

(iii) 当$p<0$时, 显然有$\underset{u\rightarrow +\infty}{\lim}\frac{1}{\left( 2-\frac{1}{u} \right) ^{\frac{1}{p}}u^{2-\frac{1}{p}}}=0$. 令$g\left( u \right) =\left( 2-\frac{1}{u} \right) ^{\frac{1}{p}}u^{2-\frac{1}{p}}$, 则
\begin{align*}
g'\left( u \right) =\frac{2}{p}u^{-\frac{1}{p}}\left( 2-\frac{1}{u} \right) ^{\frac{1}{p}-1}+\left( 2-\frac{1}{p} \right) \left( 2-\frac{1}{u} \right) ^{\frac{1}{p}}u^{1-\frac{1}{p}}>0,\forall u\in \left[ 1,+\infty \right).
\end{align*}
因此$g\left( u \right)$单调递增, 于是$\frac{1}{\left( 2-\frac{1}{u} \right) ^{\frac{1}{p}}u^{2-\frac{1}{p}}}=\frac{1}{g\left( u \right)}$单调递减趋于$0$. 又显然有$\int_1^A{\cos x\mathrm{d}x}$关于$A$有界, 所以结合\eqref{eq:integral-transform111}式, 再由Dirichlet判别法可知$I$收敛.
\end{proof}

\begin{example}
对实数 \(p\)，讨论反常积分 \(\int_0^{\infty}{\frac{\sin \left( x+\frac{1}{x} \right)}{x^p}\mathrm{d}x}\) 的敛散性.
\end{example}
\begin{remark}
令\(u=x+\frac{1}{x}\)，则
\begin{align*}
\int_1^{\infty}{\frac{\sin \left( x+\frac{1}{x} \right)}{x^p}\mathrm{d}x}=\int_u^{\infty}{\frac{\sin u}{\left( u+\sqrt{u^2-4} \right) ^p}\left( 1+\frac{u}{\sqrt{u^2-4}} \right) \mathrm{d}u}。
\end{align*}
显然\(\int_0^A{\sin u\mathrm{d}u}\)关于\(A\)有界。再证明\(\frac{1+\frac{u}{\sqrt{u^2-4}}}{\left( u+\sqrt{u^2-4} \right) ^p}\)单调递减趋于0，就能利用Dirichlet判别法得到\(\int_1^{\infty}{\frac{\sin \left( x+\frac{1}{x} \right)}{x^p}\mathrm{d}x}\)收敛。再同理讨论\(\int_0^1{\frac{\sin \left( x+\frac{1}{x} \right)}{x^p}\mathrm{d}x}\)即可。这种方法虽然能做，但是比较繁琐，不适合考场中使用。
\end{remark}
\begin{proof}
显然\(\int_0^{\infty}{\frac{\sin \left( x+\frac{1}{x} \right)}{x^p}\mathrm{d}x}\)有两个奇点\(x=0,+\infty\)。

(1) 当\(p\leqslant 0\)时，考虑区间\(\left[ 2n\pi +\frac{\pi}{4},2n\pi +\frac{3\pi}{4} \right]\)，则
\begin{align*}
x+\frac{1}{x}\in \left[ 2n\pi +\frac{\pi}{4}+\frac{1}{2n\pi +\frac{\pi}{4}},2n\pi +\frac{3\pi}{4}+\frac{1}{2n\pi +\frac{3\pi}{4}} \right]。
\end{align*}
于是当\(n>10\)时，我们有
\begin{align*}
\int_{2n\pi +\frac{\pi}{4}}^{2n\pi +\frac{3\pi}{4}}{\frac{\sin \left( x+\frac{1}{x} \right)}{x^p}\mathrm{d}x}&\geqslant \int_{2n\pi +\frac{\pi}{4}}^{2n\pi +\frac{3\pi}{4}}{\sin \left( x+\frac{1}{x} \right) \mathrm{d}x}\\
&\geqslant \int_{2n\pi +\frac{\pi}{4}}^{2n\pi +\frac{3\pi}{4}}{\sin \left( 2n\pi +\frac{3\pi}{4}+\frac{1}{2n\pi +\frac{3\pi}{4}} \right) \mathrm{d}x}\\
&=\frac{\pi}{2}\sin \left( \frac{3\pi}{4}+\frac{1}{2n\pi +\frac{3\pi}{4}} \right) >0。
\end{align*}
因此由Cauchy收敛准则可知\(\int_1^{\infty}{\frac{\sin \left( x+\frac{1}{x} \right)}{x^p}\mathrm{d}x}\)发散。故此时\(\int_0^{\infty}{\frac{\sin \left( x+\frac{1}{x} \right)}{x^p}\mathrm{d}x}\)发散。

(2) 当\(p>0\)时，先考虑\(\int_1^{\infty}{\frac{\sin \left( x+\frac{1}{x} \right)}{x^p}\mathrm{d}x}\)。

(i) 若\(p>1\)，则
\begin{align*}
\int_1^{\infty}{\left| \frac{\sin \left( x+\frac{1}{x} \right)}{x^p} \right|\mathrm{d}x}\leqslant \int_1^{\infty}{\frac{1}{x^p}\mathrm{d}x}<\infty。
\end{align*}
因此\(\int_1^{\infty}{\frac{\sin \left( x+\frac{1}{x} \right)}{x^p}\mathrm{d}x}\)绝对收敛。

(ii) 若\(p\in \left( 0,1 \right]\)，则
\begin{align}
\int_1^{\infty}{\frac{\sin \left( x+\frac{1}{x} \right)}{x^p}\mathrm{d}x}=\int_1^{\infty}{\sin x\frac{\cos \frac{1}{x}}{x^p}\mathrm{d}x}+\int_1^{\infty}{\cos x\frac{\sin \frac{1}{x}}{x^p}\mathrm{d}x}。\label{100.40}
\end{align}
显然\(\int_1^A{\cos x\mathrm{d}x}\)关于\(A\)有界，并且\(\frac{\sin \frac{1}{x}}{x^p}\)在\(\left[ 1,+\infty \right)\)上单调递减趋于0，故由Dirichlet判别法可知\(\int_1^{\infty}{\cos x\frac{\sin \frac{1}{x}}{x^p}\mathrm{d}x}\)收敛。令\(f\left( u \right) =u^p\cos u\)，则当\(u\in \left( 0,\frac{4p}{\pi} \right)\)时，有
\begin{align*}
f'\left( u \right) =pu^{p-1}\cos u-u^p\sin u=u^{p-1}\cos u\left( p-u\tan u \right) >0。
\end{align*}
于是\(f\left( u \right)\)在\(\left( 0,\frac{4p}{\pi} \right)\)上单调递增，从而\(\frac{\cos \frac{1}{x}}{x^p}=f\left( \frac{1}{x} \right)\)在\(\left( \frac{\pi}{4}p,+\infty \right)\)上单调递减趋于0。又显然\(\int_{\frac{\pi}{4}p}^A{\sin x\mathrm{d}x}\)关于\(A\)有界，故由Dirichlet判别法可知\(\int_{\frac{\pi}{4}p}^{\infty}{\sin x\frac{\cos \frac{1}{x}}{x^p}\mathrm{d}x}\)收敛，又\(\frac{\pi}{4}p<1\)，故此时\(\int_1^{\infty}{\sin x\frac{\cos \frac{1}{x}}{x^p}\mathrm{d}x}\)收敛。因此再由\eqref{100.40}式可知\(\int_1^{\infty}{\frac{\sin \left( x+\frac{1}{x} \right)}{x^p}\mathrm{d}x}\)收敛。

注意到
\begin{align}
\int_1^{\infty}{\frac{\left| \sin \left( x+\frac{1}{x} \right) \right|}{x^p}\mathrm{d}x}&\geqslant \int_1^{\infty}{\frac{\sin ^2\left( x+\frac{1}{x} \right)}{x^p}\mathrm{d}x}=\frac{1}{2}\int_1^{\infty}{\frac{1}{x^p}\mathrm{d}x}+\frac{1}{2}\int_1^{\infty}{\frac{\cos \left( 2x+\frac{2}{x} \right)}{x^p}\mathrm{d}x}\nonumber
\\
&=\frac{1}{2}\int_1^{\infty}{\frac{1}{x^p}\mathrm{d}x}+\frac{1}{2}\int_1^{\infty}{\cos \left( 2x \right) \frac{\cos \frac{2}{x}}{x^p}\mathrm{d}x}+\frac{1}{2}\int_1^{\infty}{\sin \left( 2x \right) \frac{\sin \frac{2}{x}}{x^p}\mathrm{d}x}.\label{eq:::IOIOWIOIWIIWIOIEJ3344}
\end{align}
由之前的证明,同理利用Dirichlet判别法知$\int_1^{\infty}{\cos \left( 2x \right) \frac{\cos \frac{2}{x}}{x^p}\mathrm{d}x}$和$\int_1^{\infty}{\sin \left( 2x \right) \frac{\sin \frac{2}{x}}{x^p}\mathrm{d}x}$都收敛.又因为\(\int_1^{\infty}{\frac{1}{x^p}\mathrm{d}x}\)发散,所以再由\eqref{eq:::IOIOWIOIWIIWIOIEJ3344}式知\(\int_1^{\infty}{\frac{\sin \left( x+\frac{1}{x} \right)}{x^p}\mathrm{d}x}\)条件收敛，但不绝对收敛。

再考虑\(\int_0^1{\frac{\sin \left( x+\frac{1}{x} \right)}{x^p}\mathrm{d}x}\)。

\one 若\(p\in \left( 0,1 \right)\)，则
\begin{align*}
\int_0^1{\frac{\left| \sin \left( x+\frac{1}{x} \right) \right|}{x^p}\mathrm{d}x}\leqslant \int_0^1{\frac{1}{x^p}\mathrm{d}x}<\infty。
\end{align*}
故此时\(\int_0^1{\frac{\sin \left( x+\frac{1}{x} \right)}{x^p}\mathrm{d}x}\)绝对收敛。

\two 若\(p\geqslant 1\)，则
\begin{align*}
\int_0^1{\frac{\sin \left( x+\frac{1}{x} \right)}{x^p}\mathrm{d}x}\xlongequal{x=\frac{1}{t}}\int_1^{\infty}{\frac{\sin \left( t+\frac{1}{t} \right)}{t^{2-p}}\mathrm{d}t}。
\end{align*}
此时\(2-p\leqslant 1\)。于是当\(2-p\leqslant 0\)即\(p\geqslant 2\)时，由(1)可知\(\int_0^1{\frac{\sin \left( x+\frac{1}{x} \right)}{x^p}\mathrm{d}x}\)发散。当\(2-p\in \left( 0,1 \right]\)即\(p\in \left[ 1,2 \right)\)时，由(i)可知\(\int_0^1{\frac{\sin \left( x+\frac{1}{x} \right)}{x^p}\mathrm{d}x}\)条件收敛，但不绝对收敛。

综上，当\(p\leqslant 0\)时，\(\int_0^{\infty}{\frac{\sin \left( x+\frac{1}{x} \right)}{x^p}\mathrm{d}x}\)发散；当\(p\in \left( 0,2 \right)\)时，\(\int_0^{\infty}{\frac{\sin \left( x+\frac{1}{x} \right)}{x^p}\mathrm{d}x}\)条件收敛；当\(p\geqslant 2\)时，\(\int_0^{\infty}{\frac{\sin \left( x+\frac{1}{x} \right)}{x^p}\mathrm{d}x}\)发散。
\end{proof}

\begin{example}
判断广义积分\(\int_{1}^{\infty}\frac{1}{x}e^{\cos x}\cos(2\sin x)\mathrm{d}x\)，\(\int_{0}^{\infty}\frac{1}{x}e^{\cos x}\sin(\sin x)\mathrm{d}x\)的敛散性。 
\end{example}
\begin{proof}
(1) 由于\(e^{\cos x}\sin \left( 2\sin x \right)\)是周期为\(2\pi\)的奇函数，故
\begin{align*}
\int_0^{2\pi}{e^{\cos x}\sin \left( 2\sin x \right) \mathrm{d}x}=\int_{-\pi}^{\pi}{e^{\cos x}\sin \left( 2\sin x \right) \mathrm{d}x}=0。
\end{align*}
\begin{align*}
\int_0^{2\pi}{\left| e^{\cos x}\sin \left( 2\sin x \right) \right|\mathrm{d}x}\leqslant \int_0^{2\pi}{e\mathrm{d}x}=2\pi e。
\end{align*}
于是
\begin{align*}
\int_0^A{e^{\cos x}\sin \left( 2\sin x \right) \mathrm{d}x}&=\int_0^{2\pi \left[ \frac{A}{2\pi} \right]}{e^{\cos x}\sin \left( 2\sin x \right) \mathrm{d}x}+\int_{2\pi \left[ \frac{A}{2\pi} \right]}^A{e^{\cos x}\sin \left( 2\sin x \right) \mathrm{d}x} \\
&\leqslant 0+\int_{2\pi \left[ \frac{A}{2\pi} \right]}^A{\left| e^{\cos x}\sin \left( 2\sin x \right) \right|\mathrm{d}x} 
\leqslant \int_{2\pi \left[ \frac{A}{2\pi} \right]}^{2\pi \left( \left[ \frac{A}{2\pi} \right] +1 \right)}{\left| e^{\cos x}\sin \left( 2\sin x \right) \right|\mathrm{d}x} \\
&=\int_0^{2\pi}{\left| e^{\cos x}\sin \left( 2\sin x \right) \right|\mathrm{d}x}\leqslant 2\pi e,\forall A>2\pi。
\end{align*}
又显然有\(\frac{1}{x}\)单调趋于\(0\)，故由Dirichlet判别法可知\(\int_0^{\infty}{e^{\cos x}\sin \left( 2\sin x \right) \mathrm{d}x}\)收敛。

(2) 对\(\forall n\in \mathbb{N}\)，我们有
\begin{align*}
\int_{2n\pi}^{2\left( n+2 \right) \pi}{\frac{1}{x}e^{\cos x}\cos \left( 2\sin x \right) \mathrm{d}x}\geqslant \frac{C}{n}，
\end{align*}
其中\(C\)为某一常数。(这里需要对上述积分进行数值估计，\(C\)需要具体确定出来，太麻烦暂时省略) 于是
\begin{align*}
\int_1^{\infty}{\frac{1}{x}e^{\cos x}\cos(2\sin x)\mathrm{d}x}&=\sum_{n=1}^{\infty}{\int_{2n\pi}^{2\left( n+2 \right) \pi}{\frac{1}{x}e^{\cos x}\cos \left( 2\sin x \right) \mathrm{d}x}} 
\geqslant \sum_{n=1}^{\infty}{\frac{C}{n}}=\infty。
\end{align*}
故\(\int_1^{\infty}{\frac{1}{x}e^{\cos x}\cos(2\sin x)\mathrm{d}x}\)发散。
\end{proof}

\begin{example}
设 $f(x)\in C^{1}[1,+\infty)$,$0\leqslant  f(x)\leqslant  x^{2}\ln x$，$f'(x)>0$，证明：$\int_{1}^{+\infty}\frac{1}{f'(x)}\,\mathrm{d}x$ 发散.
\end{example}
\begin{note}
首先形式计算一下，假如 $f(x)=x^{2}\ln x$，则 $f'(x)=2x\ln x+x$，量级是 $x\ln x$，代入进去刚刚好积分是发散的，可以把这个视为取等条件，然后对着这个取等，使用柯西不等式（目标是去掉难以处理的分母）.
\end{note}
\begin{proof}
对任意充分大的$b>a$,令$A=e^a,B=e^b$,则由Cauchy不等式有
\begin{align*}
\int_A^B{\frac{f' \left( x \right)}{x^2\ln ^2x}\mathrm{d}x}\int_A^B{\frac{1}{f' \left( x \right)}\mathrm{d}x}\geqslant \left( \int_A^B{\frac{1}{x\ln x}\mathrm{d}x} \right) ^2=\left( \ln\ln B-\ln\ln A \right) ^2=\left( \ln \frac{\ln B}{\ln A} \right) ^2.
\end{align*}
注意到
\begin{align*}
\int_A^B{\frac{f' \left( x \right)}{x^2\ln ^2x}\mathrm{d}x}=\int_A^B{\frac{1}{x^2\ln ^2x}\mathrm{d}f\left( x \right)}=\frac{f\left( B \right)}{B^2\ln ^2B}-\frac{f\left( A \right)}{A^2\ln ^2A}+2\int_A^B{\frac{f\left( x \right) \left( \ln x+1 \right)}{x^3\ln ^3x}\mathrm{d}x},
\end{align*}
故
\begin{align*}
\left( \ln \frac{\ln B}{\ln A} \right) ^2&\leqslant \int_A^B{\frac{1}{f' \left( x \right)}\mathrm{d}x}\left[ \frac{f\left( B \right)}{B^2\ln ^2B}-\frac{f\left( A \right)}{A^2\ln ^2A}+2\int_A^B{\frac{f\left( x \right) \left( \ln x+1 \right)}{x^3\ln ^3x}\mathrm{d}x} \right] \\
&\leqslant \int_A^B{\frac{1}{f' \left( x \right)}\mathrm{d}x}\left[ \frac{f\left( B \right)}{B^2\ln ^2B}-\frac{f\left( A \right)}{A^2\ln ^2A}+4\int_A^B{\frac{f\left( x \right)}{x^3\ln ^2x}\mathrm{d}x} \right] \\
&\leqslant \int_A^B{\frac{1}{f' \left( x \right)}\mathrm{d}x}\left( \frac{1}{\ln B}+4\int_A^B{\frac{1}{x\ln x}\mathrm{d}x} \right) \\
&=\int_A^B{\frac{1}{f' \left( x \right)}\mathrm{d}x}\left( \frac{1}{\ln B}+4\ln \frac{\ln B}{\ln A} \right) .
\end{align*}
从而
\begin{align*}
\left( \ln \frac{b}{a} \right) ^2\leqslant \int_{e^a}^{e^b}{\frac{1}{f' \left( x \right)}\mathrm{d}x}\left( \frac{1}{b}+4\ln \frac{b}{a} \right)
\Rightarrow \int_{e^a}^{e^b}{\frac{1}{f' \left( x \right)}\mathrm{d}x}\geqslant \frac{\left( \ln \frac{b}{a} \right) ^2}{\left( \frac{1}{b}+4\ln \frac{b}{a} \right)}.
\end{align*}
于是对任意充分大的$a$,取$b=2a$,则
\begin{align*}
\int_{e^a}^{e^{2a}}{\frac{1}{f' \left( x \right)}\mathrm{d}x}\geqslant \frac{\left( \ln 2 \right) ^2}{\left( \frac{1}{2a}+4\ln 2 \right)}\rightarrow \frac{\ln 2}{4},a\rightarrow +\infty .
\end{align*}
因此由Cauchy收敛准则可知$\int_1^{+\infty}{\frac{1}{f' \left( x \right)}\mathrm{d}x}$发散.
\end{proof}

\begin{example}
设 $f(x)$ 在 $[1,+\infty)$ 单调递减趋于零，$p>1$，若 $\int_{1}^{\infty}\frac{f^{p-1}(x)}{x^{\frac{1}{p}}}\,\mathrm{d}x$ 收敛，证明：$\int_{1}^{\infty}f^{p}(x)\,\mathrm{d}x$ 收敛.
\end{example}
\begin{note}
首先要搞清楚一个误区: 一定不存在$C>0$,使得
\begin{align}\label{eq:103.89tu323f8dg34t392}
\int_1^{\infty}{f^p\left( x \right)}\mathrm{d}x\leqslant C\int_1^{\infty}{\frac{f^{p-1}\left( x \right)}{x^{\frac{1}{p}}}\mathrm{d}x}
\end{align}
成立. 因为如果上式成立, 则对$\forall k>0$, 用$kf\left( x \right)$代替$f\left( x \right)$就有
\begin{align*}
k^p\int_1^{\infty}{f^p\left( x \right)}\mathrm{d}x&\leqslant Ck^{p-1}\int_1^{\infty}{\frac{f^{p-1}\left( x \right)}{x^{\frac{1}{p}}}\mathrm{d}x}\\
\Rightarrow \frac{k}{C}\int_1^{\infty}{f^p\left( x \right)}\mathrm{d}x&\leqslant \int_1^{\infty}{\frac{f^{p-1}\left( x \right)}{x^{\frac{1}{p}}}\mathrm{d}x}.
\end{align*}
令$k\rightarrow \infty$得$\int_1^{\infty}{\frac{f^{p-1}\left( x \right)}{x^{\frac{1}{p}}}\mathrm{d}x}>+\infty$, 矛盾!
因此, 只有\eqref{eq:103.89tu323f8dg34t392}式左右$f$的次数相同(齐次不等式), 才可能存在上述的$C$.
由此得到启发, 我们可以尝试建立如下不等式
\begin{align*}
\int_0^{\infty}{f^p\left( x \right)}\mathrm{d}x\leqslant C\left( \int_0^{\infty}{\frac{f^{p-1}\left( x \right)}{x^{\frac{1}{p}}}\mathrm{d}x} \right) ^{\frac{p}{p-1}}.
\end{align*}
因为$\int_0^1{\frac{1}{x^{\frac{1}{p}}}\mathrm{d}x}$收敛, 所以
\begin{align*}
\int_0^1{\frac{f^{p-1}\left( x \right)}{x^{\frac{1}{p}}}\mathrm{d}x}\leqslant \int_0^1{\frac{f\left( 1 \right)}{x^{\frac{1}{p}}}\mathrm{d}x}<+\infty .
\end{align*}
故$\int_0^{\infty}{\frac{f^{p-1}\left( x \right)}{x^{\frac{1}{p}}}\mathrm{d}x}$收敛. 从而可以不妨将积分下限改成$0$, 方便后续计算.
定义$F\left( x \right) \triangleq \begin{cases}
	f\left( 1 \right) ,x\in \left[ 0,1 \right]\\
	f\left( x \right) ,x>1\\
\end{cases}$, 则$F\left( x \right) $递减, $F\left( x \right) =f\left( x \right) $, $\forall x\geqslant 1$, 并且
\begin{align*}
\int_0^{\infty}{\frac{F^{p-1}\left( x \right)}{x^{\frac{1}{p}}}\mathrm{d}x}=\int_0^1{\frac{f\left( 1 \right)}{x^{\frac{1}{p}}}\mathrm{d}x}+\int_1^{\infty}{\frac{f^{p-1}\left( x \right)}{x^{\frac{1}{p}}}\mathrm{d}x}<+\infty .
\end{align*}
待定$C>0$, 令$g\left( x \right) =\int_0^x{F^p\left( t \right)}\mathrm{d}t-C\left( \int_0^x{\frac{F^{p-1}\left( t \right)}{t^{\frac{1}{p}}}\mathrm{d}t} \right) ^{\frac{p}{p-1}}$, 则$g\left( 0 \right) =0$, 形式计算($f$不一定连续, $g$不一定可导)可得
\begin{align*}
g' \left( x \right) &=F^p\left( x \right) -\frac{Cp}{p-1}\left( \int_0^x{\frac{F^{p-1}\left( t \right)}{t^{\frac{1}{p}}}\mathrm{d}t} \right) ^{\frac{1}{p-1}}\frac{F^{p-1}\left( x \right)}{x^{\frac{1}{p}}}\\
&=\frac{F^{p-1}\left( x \right)}{x^{\frac{1}{p}}}\left[ F\left( x \right) x^{\frac{1}{p}}-\frac{Cp}{p-1}\left( \int_0^x{\frac{F^{p-1}\left( t \right)}{t^{\frac{1}{p}}}\mathrm{d}t} \right) ^{\frac{1}{p-1}} \right] .
\end{align*}
由$F$递减可得
\begin{align*}
\frac{Cp}{p-1}\left( \int_0^x{\frac{F^{p-1}\left( t \right)}{t^{\frac{1}{p}}}\mathrm{d}t} \right) ^{\frac{1}{p-1}}\geqslant \frac{Cp}{p-1}\left( F^{p-1}\left( x \right) \int_0^x{\frac{1}{t^{\frac{1}{p}}}\mathrm{d}t} \right) ^{\frac{1}{p-1}}=C\left( \frac{p}{p-1} \right) ^{\frac{p}{p-1}}\cdot F\left( x \right) x^{\frac{1}{p}}.
\end{align*}
从而取$C=\left( 1-\frac{1}{p} \right) ^{\frac{p}{p-1}}$, 则上式可化为
\begin{align}
\left( 1-\frac{1}{p} \right) ^{\frac{1}{p-1}}\left( \int_0^x{\frac{F^{p-1}\left( t \right)}{t^{\frac{1}{p}}}\mathrm{d}t} \right) ^{\frac{1}{p-1}}\geqslant \left( 1-\frac{1}{p} \right) ^{\frac{p}{p-1}}\left( \frac{p}{p-1} \right) ^{\frac{p}{p-1}}\cdot F\left( x \right) x^{\frac{1}{p}}=F\left( x \right) x^{\frac{1}{p}},\forall x\geqslant 1.\label{eq:132312}
\end{align}
于是$g' \left( x \right) \leqslant 0$, $\forall x\geqslant 1$再结合$g\left( 0 \right) =0$就有
\begin{align*}
\int_0^x{F^p\left( t \right)}\mathrm{d}t-\left( 1-\frac{1}{p} \right) ^{\frac{1}{p-1}}\left( \int_0^x{\frac{F^{p-1}\left( t \right)}{t^{\frac{1}{p}}}\mathrm{d}t} \right) ^{\frac{p}{p-1}}=g\left( x \right) \leqslant g\left( 0 \right) =0,\forall x\geqslant 1.
\end{align*}
令$x\rightarrow \infty$得
\begin{align*}
\int_0^{\infty}{F^p\left( x \right)}\mathrm{d}x\leqslant \left( 1-\frac{1}{p} \right) ^{\frac{1}{p-1}}\left( \int_0^{\infty}{\frac{F^{p-1}\left( x \right)}{x^{\frac{1}{p}}}\mathrm{d}x} \right) ^{\frac{p}{p-1}}<+\infty .
\end{align*}
但是注意上述$g$不一定可导,所以还是需要通过定性地放缩得到严谨的证明,只需注意到\eqref{eq:132312}式始终成立.

当然也可以通过逼近方法,构造一个折线函数$h(x)$逼近$F(x)$,此时$h(x)$连续,从而用$h(x)$代替$g(x)$中的$F(x)$得到的新的$G(x)$是可导的.就能按照上述方法进行严谨证明.(逼近得到的不等式系数往往更加精确)
\end{note}
\begin{proof}
定义$F\left( x \right) \triangleq \begin{cases}
f\left( 1 \right) ,x\in \left[ 0,1 \right]\\
f\left( x \right) ,x>1\\
\end{cases}$, 则$F\left( x \right) $递减, $F\left( x \right) =f\left( x \right) $, $\forall x\geqslant 1$, 并且
\begin{align*}
\int_0^{\infty}{\frac{F^{p-1}\left( x \right)}{x^{\frac{1}{p}}}\mathrm{d}x}=\int_0^1{\frac{f\left( 1 \right)}{x^{\frac{1}{p}}}\mathrm{d}x}+\int_1^{\infty}{\frac{f^{p-1}\left( x \right)}{x^{\frac{1}{p}}}\mathrm{d}x}<+\infty .
\end{align*}
待定$C>0$, 令$g\left( x \right) =\int_0^x{F^p\left( t \right)}\mathrm{d}t-C\left( \int_0^x{\frac{F^{p-1}\left( t \right)}{t^{\frac{1}{p}}}\mathrm{d}t} \right) ^{\frac{p}{p-1}}$, 则由$F$递减可得
\begin{align*}
\frac{Cp}{p-1}\left( \int_0^x{\frac{F^{p-1}\left( t \right)}{t^{\frac{1}{p}}}\mathrm{d}t} \right) ^{\frac{1}{p-1}}\geqslant \frac{Cp}{p-1}\left( F^{p-1}\left( x \right) \int_0^x{\frac{1}{t^{\frac{1}{p}}}\mathrm{d}t} \right) ^{\frac{1}{p-1}}=C\left( \frac{p}{p-1} \right) ^{\frac{p}{p-1}}\cdot F\left( x \right) x^{\frac{1}{p}}.
\end{align*}
从而取$C=\left( 1-\frac{1}{p} \right) ^{\frac{p}{p-1}}$, 则上式可化为
\begin{align*}
\left( 1-\frac{1}{p} \right) ^{\frac{1}{p-1}}\left( \int_0^x{\frac{F^{p-1}\left( t \right)}{t^{\frac{1}{p}}}\mathrm{d}t} \right) ^{\frac{1}{p-1}}\geqslant \left( 1-\frac{1}{p} \right) ^{\frac{p}{p-1}}\left( \frac{p}{p-1} \right) ^{\frac{p}{p-1}}\cdot F\left( x \right) x^{\frac{1}{p}}=F\left( x \right) x^{\frac{1}{p}},\forall x\geqslant 1.
\end{align*}
由$\int_0^{\infty}{\frac{F^{p-1}\left( t \right)}{t^{\frac{1}{p}}}\mathrm{d}t}$收敛可知, 存在$C>0$, 使得
\begin{align*}
F\left( x \right) x^{\frac{1}{p}}\leqslant \left( 1-\frac{1}{p} \right) ^{\frac{1}{p-1}}\left( \int_0^x{\frac{F^{p-1}\left( t \right)}{t^{\frac{1}{p}}}\mathrm{d}t} \right) ^{\frac{1}{p-1}}\leqslant \left( 1-\frac{1}{p} \right) ^{\frac{1}{p-1}}\left( \int_0^{\infty}{\frac{F^{p-1}\left( t \right)}{t^{\frac{1}{p}}}\mathrm{d}t} \right) ^{\frac{1}{p-1}}<C,\forall x\geqslant 1.
\end{align*}
于是
\begin{align*}
\int_0^{\infty}{F^p\left( x \right)}\mathrm{d}x=\int_0^{\infty}{\frac{F^{p-1}\left( x \right)}{x^{\frac{1}{p}}}\cdot F\left( x \right) x^{\frac{1}{p}}\mathrm{d}x}\leqslant C\int_0^{\infty}{\frac{F^{p-1}\left( x \right)}{x^{\frac{1}{p}}}\mathrm{d}x}<+\infty .
\end{align*}
\end{proof}

\begin{proposition}\label{proposition:A-D判别法是积分收敛的“充要条件”}
设 $f(x)$ 在 $[0,+\infty)$ 中连续，证明：广义积分 $\int_{0}^{+\infty}f(x)\,\mathrm{d}x$ 收敛的充要条件有
\begin{enumerate}
\item 存在 $u(x),v(x)$ 使得 $f(x)=u(x)v(x)$，其中 $u(x)$ 单调趋于零，$\int_{0}^{A}v(x)\,\mathrm{d}x$ 有界.

\item 存在 $u(x),v(x)$ 使得 $f(x)=u(x)v(x)$，其中 $u(x)$ 单调有界，$\int_{0}^{+\infty}v(x)\,\mathrm{d}x$ 收敛，
\end{enumerate}
\end{proposition}
\begin{note}
这个命题说明:A-D判别法“几乎”是充要条件(只有确定$f$的分解逆命题才成立)，并且“逆命题”当中，依然是Dirichlet判别法强于 Abel 判别法.级数版本见\refpro{proposition:A-D判别法是级数收敛的“充要条件”}.
\end{note}
\begin{proof}
充分性由A-D判别法立得.下证明必要性.
\begin{enumerate}
\item 由$\int_0^{+\infty} f(x)\mathrm{d}x$收敛及Cauchy收敛准则可知，对$\forall \varepsilon >0$, $\exists M>0$, $\forall B>A>M$，有
\begin{align*}
\left| \int_A^B f(x)\mathrm{d}x \right| < \varepsilon.
\end{align*}
对$\forall n\in \mathbb{N}$，取$\varepsilon = \frac{1}{n^3}$，则存在$M_n>0$，对$\forall B>M_n$，有
\begin{align}
\left| \int_{M_n}^B f(x)\mathrm{d}x \right| < \frac{1}{n^3}. \label{eq:103.89tu323f8dg34t394}
\end{align}
取$\varepsilon = \frac{1}{(n+1)^3}$，则存在$M_{n+1}>M_n+1$，对$\forall B>M_{n+1}$，有
\begin{align*}
\left| \int_{M_{n+1}}^B f(x)\mathrm{d}x \right| < \frac{1}{(n+1)^3}.
\end{align*}
由$M_{n+1}>M_n+1$及\eqref{eq:103.89tu323f8dg34t394}式可知
\begin{align}\label{eq:103.89tu323f8dg34t395}
\left| \int_{M_n}^{M_{n+1}} f(x)\mathrm{d}x \right| < \frac{1}{n^3}.
\end{align}
令$u(x) \triangleq \begin{cases} 
\frac{1}{n}, & x\in [M_n, M_{n+1}),n=1,2,\cdots \\
1, & x\in [0, M_1)
\end{cases}$, $v(x) \triangleq \frac{f(x)}{u(x)}$，则$u(x)$单调递减，且$\lim\limits_{x\rightarrow +\infty} u(x) = 0$.
对$\forall A>0$，存在$n\in \mathbb{N}$，使得$A\in [M_n, M_{n+1})$. 从而由\eqref{eq:103.89tu323f8dg34t394}\eqref{eq:103.89tu323f8dg34t395}式可得
\begin{align*}
\left| \int_0^A \frac{f(x)}{u(x)}\mathrm{d}x \right| &= \left| \int_0^{M_1} \frac{f(x)}{u(x)}\mathrm{d}x + \int_{M_1}^{M_2} \frac{f(x)}{u(x)}\mathrm{d}x + \cdots + \int_{M_{n-1}}^{M_n} \frac{f(x)}{u(x)}\mathrm{d}x + \int_{M_n}^A \frac{f(x)}{u(x)}\mathrm{d}x \right| \\
&\leqslant \left| \int_0^{M_1} f(x)\mathrm{d}x \right| + \left| \int_{M_1}^{M_2} f(x)\mathrm{d}x \right| + \cdots + \left| \int_{M_{n-1}}^{M_n} (n-1)f(x)\mathrm{d}x \right| + \left| \int_{M_n}^A nf(x)\mathrm{d}x \right| \\
&\leqslant \left| \int_0^{M_1} f(x)\mathrm{d}x \right| + 1 + \cdots + \frac{1}{(n-1)^2} + \frac{1}{n^2} \\
&< \left| \int_0^{M_1} f(x)\mathrm{d}x \right| + \frac{\pi^2}{6} < +\infty.
\end{align*}
这就完成了证明.

\item 由第1问可知，存在$u(x)$, $v(x)$, 使得$f(x)=u(x)v(x)$, 其中$u(x)$单调趋于0, $\int_0^A v(x)\,\mathrm{d}x$有界. 令$u_1(x)=\sqrt{u(x)}$, $v_1(x)=\sqrt{u(x)}v(x)$, 则$f(x)=u_1(x)v_1(x)$. 由$u(x)$单调趋于0可知, $u_1(x)$单调有界. 因为$\sqrt{u(x)}$单调趋于0, $\int_0^A v(x)\mathrm{d}x$有界, 所以由第1问可知
\begin{align*}
\int_0^{\infty} v_1(x)\mathrm{d}x=\int_0^{\infty} \sqrt{u(x)}v(x)\mathrm{d}x<+\infty.
\end{align*}
故$u_1(x)$, $v_1(x)$就是第2问中我们要找的分解.
\end{enumerate}
\end{proof}









\end{document}